\chapter{INTRODUCTION}

Introductory paragraph for Chapter 1.

\section{Background of the Study}

Insert background of the study here.

As always, an indented new paragraph is written by just inserting a blank line between paragraphs.

There are two types of in-text citations. The first type is for referencing the paper in which the statement is taken. For example, consider the following sentences.
\begin{enumerate}
	\item Every generalized topology is generated by a monotonic function \cite{JlEx}.
	\item If a function is differentiable at $x_0\in\R$, then it is also continuous at $x_0$ \cite{BkEx}.
\end{enumerate}

The second type is for referencing the paper in which the author(s) is/are part of the statement, as seen in the following sentences.
\begin{enumerate}
	\item Further details of this result can be found in the dissertation of Ng \cite{ThEx}.
	\item Barsara et al.~\cite{bvv11} introduced the concept of the shadow graph of a graph.
\end{enumerate}

Both types of citation uses the command \texttt{\textbackslash cite\{key\}}, where \texttt{key} is the citation key of the reference material from the \texttt{references.bib} file. Note that the AMS citation style uses numbers, so if the name of the author is needed in the text, you have to type it manually.

If a specified result within the reference is cited, e.g., Theorem 4.1 of the dissertation of Ng, we refer to it like this: \cite{ThEx}*{Theorem 4.1}. Here, the command \texttt{\textbackslash cite\{key\}*\{ref\}} is used, where \texttt{ref} is the cited result.

\section{Statement of the Problem}

Insert statement of the problem here.

\section{Objectives of the Study}

Insert objectives here.

\section{Significance of the Study}

Insert significance here.

\section{Scope and Limitation}

Insert scope and limitation here.

\section{Methodology}

Insert methodology here.